% GENERATED FILE. Edit canonical lessons, then run npm run generate:books.
% canonical-source-hash: fnv1a64:bddaf4c6230c8b89
% canonical-lessons: FR-C428-fuite, FR-C428-robinet, FR-C428-lavabo, FR-C428-goutter, FR-C428-deboucher

\chapter{Une fuite dans la salle de bains}
\label{ch:fr-une-fuite-dans-la-salle-de-bains}
\hlchaptermodality{428}

\begin{chapteropening}
\textbf{By the end of this chapter:} \emph{I can report a bathroom problem: the leak, the tap, the washbasin, dripping and unblocking.}

Complete the last lesson of chapter 428: I can report a bathroom problem: the leak, the tap, the washbasin, dripping and unblocking.
\end{chapteropening}

\section[la fuite]{la fuite — the leak; the escape}

\label{lesson:FR-C428-fuite}



\begin{quote}
Before the new one: say the French for the pillowcase, then the French for the bathrobe.
\end{quote}

\subsection*{You'll want to know: la fuite}
\textbf{la fuite} — \textquotedblleft{}the leak; the escape\textquotedblright{}.

\begin{quote}
Il y a une fuite sous la douche.
\end{quote}

\begin{quote}
\emph{There is a leak under the shower.}
\end{quote}

\emph{Une fuite d'eau} — a water leak; \emph{prendre la fuite} — to run away.

\begin{cousinweb}
\emph{fuite} goes back to Vulgar Latin \emph{fugita} 'a fleeing', from Latin \emph{fugere} 'to flee' (Lewis \& Short): the water escapes.
\end{cousinweb}

\begin{grammarlens}[title={What you've built}]
One more word for telling, explaining and describing at B1.
\end{grammarlens}

\subsection*{Your turn}
\begin{itemize}
\raggedright
  \item \emph{Say it:} \emph{la fuite}
  \item \emph{Say it:} \emph{la fuite}, once more
  \item \emph{Say it:} a sentence of your own with \emph{la fuite}
  \item \emph{Recall:} say the French for the pillowcase, then the French for the bathrobe, then say \emph{la fuite} again
\end{itemize}

\begin{tcolorbox}[breakable,colback=teal!4,colframe=teal!35!black,title={Before you move on}]
What does la fuite mean? (\textquotedblleft{}The leak; the escape\textquotedblright{}.) Say it once more.
\end{tcolorbox}
\section[le robinet]{le robinet — the tap, the faucet}

\label{lesson:FR-C428-robinet}



\begin{quote}
Before the new one: say the French for the bathrobe, then the French for the leak.
\end{quote}

\subsection*{You'll want to know: le robinet}
\textbf{le robinet} — \textquotedblleft{}the tap, the faucet\textquotedblright{}.

\begin{quote}
Le robinet d'eau chaude ne ferme plus.
\end{quote}

\begin{quote}
\emph{The hot tap will not turn off any more.}
\end{quote}

\emph{Ouvrir le robinet} — to turn on the tap.

\begin{cousinweb}
\emph{robinet} is usually traced to an old pet name for a sheep: early taps were often shaped like a sheep's head. EtymDB lists the word without a source.
\end{cousinweb}

\begin{grammarlens}[title={What you've built}]
One more word for telling, explaining and describing at B1.
\end{grammarlens}

\subsection*{Your turn}
\begin{itemize}
\raggedright
  \item \emph{Say it:} \emph{le robinet}
  \item \emph{Say it:} \emph{le robinet}, once more
  \item \emph{Say it:} a sentence of your own with \emph{le robinet}
  \item \emph{Recall:} say the French for the bathrobe, then the French for the leak, then say \emph{le robinet} again
\end{itemize}

\begin{tcolorbox}[breakable,colback=teal!4,colframe=teal!35!black,title={Before you move on}]
What does le robinet mean? (\textquotedblleft{}The tap, the faucet\textquotedblright{}.) Say it once more.
\end{tcolorbox}
\section[le lavabo]{le lavabo — the washbasin}

\label{lesson:FR-C428-lavabo}



\begin{quote}
Before the new one: say the French for the leak, then the French for the tap.
\end{quote}

\subsection*{You'll want to know: le lavabo}
\textbf{le lavabo} — \textquotedblleft{}the washbasin\textquotedblright{}.

\begin{quote}
Il se lave les mains au lavabo.
\end{quote}

\begin{quote}
\emph{He washes his hands at the basin.}
\end{quote}

\emph{Le lavabo} is the basin in the bathroom.

\begin{cousinweb}
\emph{lavabo} is a learned borrowing from Latin \emph{lavābō} 'I will wash', from \emph{lavāre} 'to wash' (EtymDB), first a word in a church prayer.
\end{cousinweb}

\begin{grammarlens}[title={What you've built}]
One more word for telling, explaining and describing at B1.
\end{grammarlens}

\subsection*{Your turn}
\begin{itemize}
\raggedright
  \item \emph{Say it:} \emph{le lavabo}
  \item \emph{Say it:} \emph{le lavabo}, once more
  \item \emph{Say it:} a sentence of your own with \emph{le lavabo}
  \item \emph{Recall:} say the French for the leak, then the French for the tap, then say \emph{le lavabo} again
\end{itemize}

\begin{tcolorbox}[breakable,colback=teal!4,colframe=teal!35!black,title={Before you move on}]
What does le lavabo mean? (\textquotedblleft{}The washbasin\textquotedblright{}.) Say it once more.
\end{tcolorbox}
\section[goutter]{goutter — to drip}

\label{lesson:FR-C428-goutter}



\begin{quote}
Before the new one: say the French for the tap, then the French for the washbasin.
\end{quote}

\subsection*{You'll want to know: goutter}
\textbf{goutter} — \textquotedblleft{}to drip\textquotedblright{}.

\begin{quote}
Le robinet goutte toute la nuit.
\end{quote}

\begin{quote}
\emph{The tap drips all night.}
\end{quote}

\emph{Une goutte} — a drop.

\begin{cousinweb}
\emph{goutter} is built on \emph{goutte} 'drop', which is inherited from Latin \emph{gutta} (Lewis \& Short). EtymDB does not list the verb.
\end{cousinweb}

\begin{grammarlens}[title={What you've built}]
One more word for telling, explaining and describing at B1.
\end{grammarlens}

\subsection*{Your turn}
\begin{itemize}
\raggedright
  \item \emph{Say it:} \emph{goutter}
  \item \emph{Say it:} \emph{goutter}, once more
  \item \emph{Say it:} a sentence of your own with \emph{goutter}
  \item \emph{Recall:} say the French for the tap, then the French for the washbasin, then say \emph{goutter} again
\end{itemize}

\begin{tcolorbox}[breakable,colback=teal!4,colframe=teal!35!black,title={Before you move on}]
What does goutter mean? (\textquotedblleft{}To drip\textquotedblright{}.) Say it once more.
\end{tcolorbox}
\section[déboucher]{déboucher — to unblock; to uncork; to come out (onto)}

\label{lesson:FR-C428-deboucher}



\begin{quote}
Before the new one: say the French for the washbasin, then the French for to drip.
\end{quote}

\subsection*{You'll want to know: déboucher}
\textbf{déboucher} — \textquotedblleft{}to unblock; to uncork; to come out (onto)\textquotedblright{}.

\begin{quote}
Le plombier a débouché la douche en dix minutes.
\end{quote}

\begin{quote}
\emph{The plumber unblocked the shower in ten minutes.}
\end{quote}

\emph{Déboucher une bouteille} — to open a bottle; \emph{la rue débouche sur la place} — the street comes out onto the square.

\begin{cousinweb}
\emph{déboucher} is built from \emph{dé-} and \emph{boucher} 'to stop up' (EtymDB), which comes from Middle French \emph{bousche} 'bundle of straw', the old way to plug a hole.
\end{cousinweb}

\begin{grammarlens}[title={What you've built}]
One more word for telling, explaining and describing at B1.
\end{grammarlens}

\subsection*{Your turn}
\begin{itemize}
\raggedright
  \item \emph{Say it:} \emph{déboucher}
  \item \emph{Say it:} \emph{déboucher}, once more
  \item \emph{Say it:} a sentence of your own with \emph{déboucher}
  \item \emph{Recall:} say the French for the washbasin, then the French for to drip, then say \emph{déboucher} again
\end{itemize}

\begin{tcolorbox}[breakable,colback=teal!4,colframe=teal!35!black,title={Before you move on}]
What does déboucher mean? (\textquotedblleft{}To unblock; to uncork; to come out (onto)\textquotedblright{}.) Say it once more.
\end{tcolorbox}
